\documentclass[12pt,a4paper]{article}
\usepackage[utf8]{vietnam}
\usepackage{amsmath,amssymb}
\usepackage{tabularx}
\usepackage{tikz}
\usepackage{xcolor}
\usepackage[top=0.6cm, bottom=0.6cm, left=0.6cm, right=0.6cm]{geometry}
\usepackage[dethi]{ex_test}

\SAOPTN{widthSA=2.8,heightSA=0.62}

\setlength{\parskip}{0pt}

% Dòng kẻ chấm mờ đều đặn chuẩn viết tay (khoảng cách 6.5mm rộng rãi vừa vặn)
\newcommand{\dongke}[1]{%
  \begingroup%
  \setlength{\parskip}{0pt}%
  \baselineskip=18.5pt%
  \vspace{1.0mm}%
  \foreach \i in {1,...,#1}{%
    \noindent{\color{gray!60}\dotfill}\par%
  }%
  \vspace{0.6mm}%
  \endgroup%
}

% Định nghĩa ô bài tập tự luận trong bảng
\newcommand{\oBaiTap}[3][6]{%
  \vspace{1.0mm}%
  \noindent\textbf{#2} #3\par%
  \dongke{#1}%
}

\begin{document}
\pagestyle{empty}
% =============================================================================
% ========================= TOÁN 6 - TUẦN 07 - BUỔI 01: MÃ ĐỀ 101 =============
% ========================= MẶT TRƯỚC (TRANG 1/2) ==============================
% =============================================================================

\noindent
\begin{tabularx}{\linewidth}{|X|c|}
\hline
\textbf{PHIẾU KIỂM TRA ĐỊNH KỲ TOÁN 6 - TUẦN 07 - BUỔI 01} & \textbf{Mã đề: 101 - Trang 1/2} \\
\hline
\end{tabularx}

\vspace{1mm}

\noindent
\begin{tabularx}{\linewidth}{|p{3.2cm}|X|}
\hline
\centering\textbf{ĐIỂM SỐ} & \centering\textbf{LỜI NHẬN XÉT CỦA THẦY CÔ} \tabularnewline
\hline
\rule{0pt}{1.8cm} & \tabularnewline
\hline
\end{tabularx}

\vspace{1.2mm}

\noindent\textbf{Câu 1 (0,25 điểm).} Tập hợp $A = \{x \in \mathbb{N} \mid 12 \le x \le 45\}$ có bao nhiêu phần tử? \shortans{34}\par\vspace{0.3mm}
\noindent\textbf{Câu 2 (0,25 điểm).} Số La Mã $\text{XXIV}$ có giá trị trong hệ thập phân bằng bao nhiêu? \shortans{24}\par\vspace{0.3mm}
\noindent\textbf{Câu 3 (0,25 điểm).} Viết gọn tích $2^3 \cdot 2^5$ dưới dạng một lũy thừa. \shortans{2^8}\par\vspace{0.3mm}
\noindent\textbf{Câu 4 (0,25 điểm).} Giá trị của biểu thức $30 - 2 \cdot 3^2$ bằng bao nhiêu? \shortans{12}\par\vspace{0.3mm}
\noindent\textbf{Câu 5 (0,25 điểm).} Số dư của phép chia tổng $(45 + 120 + 23)$ cho $5$ bằng bao nhiêu? \shortans{3}\par\vspace{0.3mm}
\noindent\textbf{Câu 6 (0,25 điểm).} Tìm chữ số $x$ lớn nhất để số $\overline{23x}$ chia hết cho $3$. \shortans{7}\par\vspace{0.3mm}
\noindent\textbf{Câu 7 (0,25 điểm).} Trong các số $11; 15; 21; 27$, số nguyên tố là số nào? \shortans{11}\par\vspace{0.3mm}
\noindent\textbf{Câu 8 (0,25 điểm).} Một hình lục giác đều có bao nhiêu đường chéo chính? \shortans{3}\par\vspace{0.3mm}
\noindent\textbf{Câu 9 (0,25 điểm).} Tính chu vi hình chữ nhật có kích thước $15\text{ m} \times 8\text{ m}$. \shortans{46\text{ m}}\par\vspace{0.3mm}
\noindent\textbf{Câu 10 (0,25 điểm).} Tính diện tích hình thoi có hai đường chéo $8\text{ cm}$ và $10\text{ cm}$. \shortans{40\text{ cm}^2}\par\vspace{0.3mm}
\vspace{0.4mm}

\noindent\textbf{Câu 11 (2,0 điểm). Thực hiện phép tính (hợp lý nếu có thể):}
\vspace{0.6mm}

\noindent
\begin{tabularx}{\linewidth}{|X|X|}
\hline
\oBaiTap[10]{a)}{$145 - 2^3 \cdot 5 + 3^2 \cdot 4$} & 
\oBaiTap[10]{b)}{$48 \cdot 67 + 48 \cdot 33 - 800$} \\
\hline
\oBaiTap[10]{c)}{$180 : [45 - (18 - 15)^2]$} & 
\oBaiTap[10]{d)}{$250 : \{ 90 - [32 + (14 - 12)^3] \}$} \\
\hline
\end{tabularx}

\newpage
% =============================================================================
% ========================= TOÁN 6 - TUẦN 07 - BUỔI 01: MÃ ĐỀ 101 =============
% ========================= MẶT SAU (TRANG 2/2) ================================
% =============================================================================

\noindent
\begin{tabularx}{\linewidth}{|X|c|}
\hline
\textbf{TOÁN 6 - TUẦN 07 - BUỔI 01} & \textbf{Mã đề: 101 - Trang 2/2} \\
\hline
\end{tabularx}

\vspace{1.5mm}

\noindent\textbf{Câu 12 (2,0 điểm). Tìm số tự nhiên $x$, biết:}
\vspace{0.8mm}

\noindent
\begin{tabularx}{\linewidth}{|X|X|}
\hline
\oBaiTap[5]{a)}{$4 \cdot x + 15 = 47$} & 
\oBaiTap[5]{b)}{$85 - 3 \cdot (x + 4) = 40$} \\
\hline
\oBaiTap[5]{c)}{$3 \cdot 2^{x+1} - 12 = 36$} & 
\oBaiTap[5]{d)}{Tìm số tự nhiên $x$ lớn nhất, biết $24 \ \vdots \ x$ và $36 \ \vdots \ x$.} \\
\hline
\end{tabularx}

\vspace{1.5mm}

\noindent
\begin{tabularx}{\linewidth}{|X|}
\hline
\vspace{0.5mm}
\textbf{Câu 13 (2,5 điểm).} Bác An có một mảnh sân hình chữ nhật có chiều dài $16\text{ m}$ và chiều rộng $10\text{ m}$.\par
\textbf{a) (1,25 điểm)} Tính chu vi và diện tích của mảnh sân hình chữ nhật đó.\par
\textbf{b) (1,25 điểm)} Bác An dùng những viên gạch hình vuông có cạnh $40\text{ cm}$ để lát kín mảnh sân đó. Hỏi bác An cần bao nhiêu viên gạch (bỏ qua mép vữa)? Biết mỗi viên gạch có giá $25\ 000$ đồng, tính tổng số tiền bác An phải trả để mua gạch.\par
\vspace{0.4mm}
\noindent\textit{Bài giải:}\dongke{11}
\\ \hline
\end{tabularx}

\vspace{1.5mm}

\noindent
\begin{tabularx}{\linewidth}{|X|}
\hline
\vspace{0.5mm}
\textbf{Câu 14 (1,0 điểm).} Cho tổng $A = 1 + 3 + 3^2 + 3^3 + \dots + 3^{21}$. Chứng minh rằng $A$ chia hết cho $4$.\par
\vspace{0.4mm}
\noindent\textit{Bài giải:}\dongke{8}
\\ \hline
\end{tabularx}

\newpage

% =============================================================================
% ========================= TOÁN 6 - TUẦN 07 - BUỔI 01: MÃ ĐỀ 102 =============
% ========================= MẶT TRƯỚC (TRANG 1/2) ==============================
% =============================================================================

\noindent
\begin{tabularx}{\linewidth}{|X|c|}
\hline
\textbf{PHIẾU KIỂM TRA ĐỊNH KỲ TOÁN 6 - TUẦN 07 - BUỔI 01} & \textbf{Mã đề: 102 - Trang 1/2} \\
\hline
\end{tabularx}

\vspace{1mm}

\noindent
\begin{tabularx}{\linewidth}{|p{3.2cm}|X|}
\hline
\centering\textbf{ĐIỂM SỐ} & \centering\textbf{LỜI NHẬN XÉT CỦA THẦY CÔ} \tabularnewline
\hline
\rule{0pt}{1.8cm} & \tabularnewline
\hline
\end{tabularx}

\vspace{1.2mm}

\noindent\textbf{Câu 1 (0,25 điểm).} Tập hợp $A = \{x \in \mathbb{N} \mid 15 \le x \le 58\}$ có bao nhiêu phần tử? \shortans{44}\par\vspace{0.3mm}
\noindent\textbf{Câu 2 (0,25 điểm).} Số La Mã $\text{XXVII}$ có giá trị trong hệ thập phân bằng bao nhiêu? \shortans{27}\par\vspace{0.3mm}
\noindent\textbf{Câu 3 (0,25 điểm).} Viết gọn tích $3^2 \cdot 3^4$ dưới dạng một lũy thừa. \shortans{3^6}\par\vspace{0.3mm}
\noindent\textbf{Câu 4 (0,25 điểm).} Giá trị của biểu thức $40 - 3 \cdot 2^3$ bằng bao nhiêu? \shortans{16}\par\vspace{0.3mm}
\noindent\textbf{Câu 5 (0,25 điểm).} Số dư của phép chia tổng $(35 + 150 + 24)$ cho $5$ bằng bao nhiêu? \shortans{4}\par\vspace{0.3mm}
\noindent\textbf{Câu 6 (0,25 điểm).} Tìm chữ số $x$ lớn nhất để số $\overline{34x}$ chia hết cho $3$. \shortans{8}\par\vspace{0.3mm}
\noindent\textbf{Câu 7 (0,25 điểm).} Trong các số $13; 25; 33; 35$, số nguyên tố là số nào? \shortans{13}\par\vspace{0.3mm}
\noindent\textbf{Câu 8 (0,25 điểm).} Một hình lục giác đều có bao nhiêu cạnh bằng nhau? \shortans{6}\par\vspace{0.3mm}
\noindent\textbf{Câu 9 (0,25 điểm).} Tính chu vi hình chữ nhật có kích thước $16\text{ m} \times 7\text{ m}$. \shortans{46\text{ m}}\par\vspace{0.3mm}
\noindent\textbf{Câu 10 (0,25 điểm).} Tính diện tích hình thoi có hai đường chéo $6\text{ cm}$ và $12\text{ cm}$. \shortans{36\text{ cm}^2}\par\vspace{0.3mm}
\vspace{0.4mm}

\noindent\textbf{Câu 11 (2,0 điểm). Thực hiện phép tính (hợp lý nếu có thể):}
\vspace{0.6mm}

\noindent
\begin{tabularx}{\linewidth}{|X|X|}
\hline
\oBaiTap[10]{a)}{$155 - 2^4 \cdot 3 + 4^2 \cdot 2$} & 
\oBaiTap[10]{b)}{$56 \cdot 72 + 56 \cdot 28 - 600$} \\
\hline
\oBaiTap[10]{c)}{$210 : [58 - (19 - 15)^2]$} & 
\oBaiTap[10]{d)}{$300 : \{ 100 - [32 + (15 - 13)^3] \}$} \\
\hline
\end{tabularx}

\newpage
% =============================================================================
% ========================= TOÁN 6 - TUẦN 07 - BUỔI 01: MÃ ĐỀ 102 =============
% ========================= MẶT SAU (TRANG 2/2) ================================
% =============================================================================

\noindent
\begin{tabularx}{\linewidth}{|X|c|}
\hline
\textbf{TOÁN 6 - TUẦN 07 - BUỔI 01} & \textbf{Mã đề: 102 - Trang 2/2} \\
\hline
\end{tabularx}

\vspace{1.5mm}

\noindent\textbf{Câu 12 (2,0 điểm). Tìm số tự nhiên $x$, biết:}
\vspace{0.8mm}

\noindent
\begin{tabularx}{\linewidth}{|X|X|}
\hline
\oBaiTap[5]{a)}{$5 \cdot x + 18 = 53$} & 
\oBaiTap[5]{b)}{$95 - 4 \cdot (x + 3) = 47$} \\
\hline
\oBaiTap[5]{c)}{$2 \cdot 3^{x+1} - 15 = 39$} & 
\oBaiTap[5]{d)}{Tìm số tự nhiên $x$ lớn nhất, biết $30 \ \vdots \ x$ và $45 \ \vdots \ x$.} \\
\hline
\end{tabularx}

\vspace{1.5mm}

\noindent
\begin{tabularx}{\linewidth}{|X|}
\hline
\vspace{0.5mm}
\textbf{Câu 13 (2,5 điểm).} Bác Bình có một mảnh sân hình chữ nhật có chiều dài $18\text{ m}$ và chiều rộng $10\text{ m}$.\par
\textbf{a) (1,25 điểm)} Tính chu vi và diện tích của mảnh sân hình chữ nhật đó.\par
\textbf{b) (1,25 điểm)} Bác Bình dùng những viên gạch hình vuông có cạnh $50\text{ cm}$ để lát kín mảnh sân đó. Hỏi bác Bình cần bao nhiêu viên gạch (bỏ qua mép vữa)? Biết mỗi viên gạch có giá $35\ 000$ đồng, tính tổng số tiền bác Bình phải trả để mua gạch.\par
\vspace{0.4mm}
\noindent\textit{Bài giải:}\dongke{11}
\\ \hline
\end{tabularx}

\vspace{1.5mm}

\noindent
\begin{tabularx}{\linewidth}{|X|}
\hline
\vspace{0.5mm}
\textbf{Câu 14 (1,0 điểm).} Cho tổng $B = 1 + 4 + 4^2 + 4^3 + \dots + 4^{21}$. Chứng minh rằng $B$ chia hết cho $5$.\par
\vspace{0.4mm}
\noindent\textit{Bài giải:}\dongke{8}
\\ \hline
\end{tabularx}

\newpage

% =============================================================================
% ========================= TOÁN 6 - TUẦN 07 - BUỔI 01: MÃ ĐỀ 103 =============
% ========================= MẶT TRƯỚC (TRANG 1/2) ==============================
% =============================================================================

\noindent
\begin{tabularx}{\linewidth}{|X|c|}
\hline
\textbf{PHIẾU KIỂM TRA ĐỊNH KỲ TOÁN 6 - TUẦN 07 - BUỔI 01} & \textbf{Mã đề: 103 - Trang 1/2} \\
\hline
\end{tabularx}

\vspace{1mm}

\noindent
\begin{tabularx}{\linewidth}{|p{3.2cm}|X|}
\hline
\centering\textbf{ĐIỂM SỐ} & \centering\textbf{LỜI NHẬN XÉT CỦA THẦY CÔ} \tabularnewline
\hline
\rule{0pt}{1.8cm} & \tabularnewline
\hline
\end{tabularx}

\vspace{1.2mm}

\noindent\textbf{Câu 1 (0,25 điểm).} Tập hợp $A = \{x \in \mathbb{N} \mid 18 \le x \le 62\}$ có bao nhiêu phần tử? \shortans{45}\par\vspace{0.3mm}
\noindent\textbf{Câu 2 (0,25 điểm).} Số La Mã $\text{XXIX}$ có giá trị trong hệ thập phân bằng bao nhiêu? \shortans{29}\par\vspace{0.3mm}
\noindent\textbf{Câu 3 (0,25 điểm).} Viết gọn tích $5^2 \cdot 5^3$ dưới dạng một lũy thừa. \shortans{5^5}\par\vspace{0.3mm}
\noindent\textbf{Câu 4 (0,25 điểm).} Giá trị của biểu thức $50 - 4 \cdot 2^3$ bằng bao nhiêu? \shortans{18}\par\vspace{0.3mm}
\noindent\textbf{Câu 5 (0,25 điểm).} Số dư của phép chia tổng $(65 + 140 + 32)$ cho $5$ bằng bao nhiêu? \shortans{2}\par\vspace{0.3mm}
\noindent\textbf{Câu 6 (0,25 điểm).} Tìm chữ số $x$ lớn nhất để số $\overline{45x}$ chia hết cho $3$. \shortans{9}\par\vspace{0.3mm}
\noindent\textbf{Câu 7 (0,25 điểm).} Trong các số $17; 27; 39; 49$, số nguyên tố là số nào? \shortans{17}\par\vspace{0.3mm}
\noindent\textbf{Câu 8 (0,25 điểm).} Một hình lục giác đều có bao nhiêu góc ở các đỉnh bằng nhau? \shortans{6}\par\vspace{0.3mm}
\noindent\textbf{Câu 9 (0,25 điểm).} Tính chu vi hình chữ nhật có kích thước $18\text{ m} \times 6\text{ m}$. \shortans{48\text{ m}}\par\vspace{0.3mm}
\noindent\textbf{Câu 10 (0,25 điểm).} Tính diện tích hình thoi có hai đường chéo $8\text{ cm}$ và $14\text{ cm}$. \shortans{56\text{ cm}^2}\par\vspace{0.3mm}
\vspace{0.4mm}

\noindent\textbf{Câu 11 (2,0 điểm). Thực hiện phép tính (hợp lý nếu có thể):}
\vspace{0.6mm}

\noindent
\begin{tabularx}{\linewidth}{|X|X|}
\hline
\oBaiTap[10]{a)}{$165 - 3^3 \cdot 2 + 5^2 \cdot 3$} & 
\oBaiTap[10]{b)}{$64 \cdot 75 + 64 \cdot 25 - 400$} \\
\hline
\oBaiTap[10]{c)}{$220 : [69 - (21 - 16)^2]$} & 
\oBaiTap[10]{d)}{$350 : \{ 110 - [32 + (16 - 14)^3] \}$} \\
\hline
\end{tabularx}

\newpage
% =============================================================================
% ========================= TOÁN 6 - TUẦN 07 - BUỔI 01: MÃ ĐỀ 103 =============
% ========================= MẶT SAU (TRANG 2/2) ================================
% =============================================================================

\noindent
\begin{tabularx}{\linewidth}{|X|c|}
\hline
\textbf{TOÁN 6 - TUẦN 07 - BUỔI 01} & \textbf{Mã đề: 103 - Trang 2/2} \\
\hline
\end{tabularx}

\vspace{1.5mm}

\noindent\textbf{Câu 12 (2,0 điểm). Tìm số tự nhiên $x$, biết:}
\vspace{0.8mm}

\noindent
\begin{tabularx}{\linewidth}{|X|X|}
\hline
\oBaiTap[5]{a)}{$6 \cdot x + 14 = 56$} & 
\oBaiTap[5]{b)}{$90 - 5 \cdot (x + 2) = 45$} \\
\hline
\oBaiTap[5]{c)}{$4 \cdot 2^{x+2} - 18 = 46$} & 
\oBaiTap[5]{d)}{Tìm số tự nhiên $x$ lớn nhất, biết $28 \ \vdots \ x$ và $42 \ \vdots \ x$.} \\
\hline
\end{tabularx}

\vspace{1.5mm}

\noindent
\begin{tabularx}{\linewidth}{|X|}
\hline
\vspace{0.5mm}
\textbf{Câu 13 (2,5 điểm).} Bác Cường có một mảnh sân hình chữ nhật có chiều dài $20\text{ m}$ và chiều rộng $12\text{ m}$.\par
\textbf{a) (1,25 điểm)} Tính chu vi và diện tích của mảnh sân hình chữ nhật đó.\par
\textbf{b) (1,25 điểm)} Bác Cường dùng những viên gạch hình vuông có cạnh $40\text{ cm}$ để lát kín mảnh sân đó. Hỏi bác Cường cần bao nhiêu viên gạch (bỏ qua mép vữa)? Biết mỗi viên gạch có giá $20\ 000$ đồng, tính tổng số tiền bác Cường phải trả để mua gạch.\par
\vspace{0.4mm}
\noindent\textit{Bài giải:}\dongke{11}
\\ \hline
\end{tabularx}

\vspace{1.5mm}

\noindent
\begin{tabularx}{\linewidth}{|X|}
\hline
\vspace{0.5mm}
\textbf{Câu 14 (1,0 điểm).} Cho tổng $C = 1 + 5 + 5^2 + 5^3 + \dots + 5^{21}$. Chứng minh rằng $C$ chia hết cho $6$.\par
\vspace{0.4mm}
\noindent\textit{Bài giải:}\dongke{8}
\\ \hline
\end{tabularx}

\newpage

% =============================================================================
% ========================= TOÁN 6 - TUẦN 07 - BUỔI 01: MÃ ĐỀ 104 =============
% ========================= MẶT TRƯỚC (TRANG 1/2) ==============================
% =============================================================================

\noindent
\begin{tabularx}{\linewidth}{|X|c|}
\hline
\textbf{PHIẾU KIỂM TRA ĐỊNH KỲ TOÁN 6 - TUẦN 07 - BUỔI 01} & \textbf{Mã đề: 104 - Trang 1/2} \\
\hline
\end{tabularx}

\vspace{1mm}

\noindent
\begin{tabularx}{\linewidth}{|p{3.2cm}|X|}
\hline
\centering\textbf{ĐIỂM SỐ} & \centering\textbf{LỜI NHẬN XÉT CỦA THẦY CÔ} \tabularnewline
\hline
\rule{0pt}{1.8cm} & \tabularnewline
\hline
\end{tabularx}

\vspace{1.2mm}

\noindent\textbf{Câu 1 (0,25 điểm).} Tập hợp $A = \{x \in \mathbb{N} \mid 14 \le x \le 69\}$ có bao nhiêu phần tử? \shortans{56}\par\vspace{0.3mm}
\noindent\textbf{Câu 2 (0,25 điểm).} Số La Mã $\text{XXVI}$ có giá trị trong hệ thập phân bằng bao nhiêu? \shortans{26}\par\vspace{0.3mm}
\noindent\textbf{Câu 3 (0,25 điểm).} Viết gọn tích $2^4 \cdot 2^3$ dưới dạng một lũy thừa. \shortans{2^7}\par\vspace{0.3mm}
\noindent\textbf{Câu 4 (0,25 điểm).} Giá trị của biểu thức $60 - 5 \cdot 2^3$ bằng bao nhiêu? \shortans{20}\par\vspace{0.3mm}
\noindent\textbf{Câu 5 (0,25 điểm).} Số dư của phép chia tổng $(75 + 160 + 41)$ cho $5$ bằng bao nhiêu? \shortans{1}\par\vspace{0.3mm}
\noindent\textbf{Câu 6 (0,25 điểm).} Tìm chữ số $x$ lớn nhất để số $\overline{16x}$ chia hết cho $3$. \shortans{8}\par\vspace{0.3mm}
\noindent\textbf{Câu 7 (0,25 điểm).} Trong các số $19; 35; 45; 51$, số nguyên tố là số nào? \shortans{19}\par\vspace{0.3mm}
\noindent\textbf{Câu 8 (0,25 điểm).} Một hình thoi có bao nhiêu cặp cạnh đối diện song song với nhau? \shortans{2}\par\vspace{0.3mm}
\noindent\textbf{Câu 9 (0,25 điểm).} Tính chu vi hình chữ nhật có kích thước $20\text{ m} \times 7\text{ m}$. \shortans{54\text{ m}}\par\vspace{0.3mm}
\noindent\textbf{Câu 10 (0,25 điểm).} Tính diện tích hình thoi có hai đường chéo $10\text{ cm}$ và $12\text{ cm}$. \shortans{60\text{ cm}^2}\par\vspace{0.3mm}
\vspace{0.4mm}

\noindent\textbf{Câu 11 (2,0 điểm). Thực hiện phép tính (hợp lý nếu có thể):}
\vspace{0.6mm}

\noindent
\begin{tabularx}{\linewidth}{|X|X|}
\hline
\oBaiTap[10]{a)}{$175 - 2^5 \cdot 2 + 3^3 \cdot 3$} & 
\oBaiTap[10]{b)}{$72 \cdot 81 + 72 \cdot 19 - 200$} \\
\hline
\oBaiTap[10]{c)}{$240 : [84 - (23 - 17)^2]$} & 
\oBaiTap[10]{d)}{$400 : \{ 120 - [32 + (17 - 15)^3] \}$} \\
\hline
\end{tabularx}

\newpage
% =============================================================================
% ========================= TOÁN 6 - TUẦN 07 - BUỔI 01: MÃ ĐỀ 104 =============
% ========================= MẶT SAU (TRANG 2/2) ================================
% =============================================================================

\noindent
\begin{tabularx}{\linewidth}{|X|c|}
\hline
\textbf{TOÁN 6 - TUẦN 07 - BUỔI 01} & \textbf{Mã đề: 104 - Trang 2/2} \\
\hline
\end{tabularx}

\vspace{1.5mm}

\noindent\textbf{Câu 12 (2,0 điểm). Tìm số tự nhiên $x$, biết:}
\vspace{0.8mm}

\noindent
\begin{tabularx}{\linewidth}{|X|X|}
\hline
\oBaiTap[5]{a)}{$7 \cdot x + 16 = 65$} & 
\oBaiTap[5]{b)}{$100 - 6 \cdot (x + 1) = 52$} \\
\hline
\oBaiTap[5]{c)}{$5 \cdot 3^{x-1} + 15 = 60$} & 
\oBaiTap[5]{d)}{Tìm số tự nhiên $x$ lớn nhất, biết $32 \ \vdots \ x$ và $48 \ \vdots \ x$.} \\
\hline
\end{tabularx}

\vspace{1.5mm}

\noindent
\begin{tabularx}{\linewidth}{|X|}
\hline
\vspace{0.5mm}
\textbf{Câu 13 (2,5 điểm).} Bác Dũng có một mảnh sân hình chữ nhật có chiều dài $25\text{ m}$ và chiều rộng $10\text{ m}$.\par
\textbf{a) (1,25 điểm)} Tính chu vi và diện tích của mảnh sân hình chữ nhật đó.\par
\textbf{b) (1,25 điểm)} Bác Dũng dùng những viên gạch hình vuông có cạnh $50\text{ cm}$ để lát kín mảnh sân đó. Hỏi bác Dũng cần bao nhiêu viên gạch (bỏ qua mép vữa)? Biết mỗi viên gạch có giá $30\ 000$ đồng, tính tổng số tiền bác Dũng phải trả để mua gạch.\par
\vspace{0.4mm}
\noindent\textit{Bài giải:}\dongke{11}
\\ \hline
\end{tabularx}

\vspace{1.5mm}

\noindent
\begin{tabularx}{\linewidth}{|X|}
\hline
\vspace{0.5mm}
\textbf{Câu 14 (1,0 điểm).} Cho tổng $D = 1 + 2 + 2^2 + 2^3 + \dots + 2^{21}$. Chứng minh rằng $D$ chia hết cho $3$.\par
\vspace{0.4mm}
\noindent\textit{Bài giải:}\dongke{8}
\\ \hline
\end{tabularx}

\newpage

\end{document}
